%&pdftex
% J. J. GUADALUPE, M. PEREZ, F. J. RUIZ and J. L. VARONA
% MANDADO a los Proc. de la A.M.S. en diciembre de 1989
% ACEPTADO en marzo de 1991
% Versi\'on definitiva tras corregir las erratas indicadas por el "referee"
% PUBLICADO EN: Proc. Amer. Math. Soc. 116 (1992), 457-464.

\font\rmdiez=cmr10 at 8pt
\def\proof{\hbox{\it Proof.\enspace}}
\font\msbm=msbm10
\def\R{\hbox{\msbm R}}
%\def\R{\hbox{\bf R}}
\def\por{\thinspace}
\magnification\magstep1

\centerline{\bf TWO NOTES ON CONVERGENCE AND DIVERGENCE A.E. OF FOURIER}  
\centerline{\bf SERIES WITH RESPECT TO SOME ORTHOGONAL SYSTEMS}
\medskip
\centerline {by}
\medskip
\centerline {J. J. GUADALUPE, M. P\'EREZ, F. J. RUIZ and J. L. VARONA}
{\parindent=0pt\footnote{}{  AMS (1985) Subject Classification: {\bf 42C10}}
\footnote{}{  Keywords and phrases: Fourier series, orthonormal
polynomials, maximal operators, $A_p$-weights}\footnote{}{  Research supported
by CAICYT PB85-0338 and DGICYT PB89-0181-C02-02}\footnote{}{  THIS PAPER HAS
BEEN PUBLISHED IN: {\it Proc.\ Amer.\ Math.\ Soc.}\ {\bf 116} (1992), 457--464}}
\bigskip
\midinsert
\narrower\narrower
\noindent ABSTRACT: {\sl We study some problems related to convergence and
divergence a.e. for Fourier series in systems $\{\phi_k\}$, where $\{\phi_k\}$
is either a system of orthonormal polynomials with respect to a measure $d\mu$
on $[-1,1]$ or a Bessel system on $[0,1]$. We obtain boundedness in weighted
$L^p$ spaces for the maximal operators associated to Fourier-Jacobi and
Fourier-Bessel series. On the other hand, we find general results about
divergence a.e. of the Fourier  series associated to Bessel systems and systems
of orthonormal polynomials on $[-1,1]$.}
\endinsert
\medskip
\noindent {\bf \S 0. Introduction.}
\medskip
\noindent Let $d\mu$ be a positive measure on a finite interval $[a,b]\subset \R$
and
$$
S_nf(x) = \sum_{k=1}^nc_k\phi_k(x) \ ,\qquad \qquad
c_k=\int_a^bf(t)\phi_k(t)\por d\mu(t)
$$
the $n$-th partial sum of the Fourier series of $f$ with respect to a real,
complete and orthonormal system $\{\phi_k\}$ on $L^2(d\mu)$.

In what follows $\{\phi_k\}$ will be either the system of orthonormal
polynomials given by the normalization of the sequence $\{1,x,x^2,\dots\}$ in
$L^2([-1,1];d\mu)$, or the Bessel system in $L^2([0,1];d\mu)$, where
$d\mu=x\por dx$.

We shall denote by $S^*f(x)=\displaystyle\sup_n|S_nf(x)|$ the maximal operator
associated to the Fourier series of $f$.We are interested in the two following
questions:


\item{(a)}Given $p\in [1,2)$\por, does $S_nf(x)$ converge to $f(x)\por \mu
\por a.e.$ for any $f$ in $L^p(d\mu)$ ?

\item{(b)}Is $S^*$ of strong or weak type (p,p) ?

Very well known standard arguments show that an affirmative answer to question
(b) leads to an affirmative answer to question (a). Several authors have
studied the former questions in the case of orthonormal polynomials on $[-1,1]$
or Bessel systems on $[0,1]$. Titchmarsh {\bf[16]} and Benedek-Panzone {\bf[3]}
analyze question (a) in the case of Fourier-Bessel series, determinating the
range of $p's$ such that (a) is true, but their techniques (equiconvergence
theorems) cannot be used to approach matter (b).

For Jacobi polynomials, equiconvergence theorems again show that there
exists a range of $p^,s$ such that (a) is true (see {\bf[15]}). Badkov studies 
in {\bf[2]} the problem of the boundedness of the operator $S^*$ for generalized
Jacobi polynomials. He obtains boundedness (even in some weighted $L^p$-spaces)
for the same range of  $p's$ in which (a) is true.

With respect to negative answers to question (a), Pollard {\bf[14]} shows that
for each $p<4/3$ there exists a function $f\in L^p(dx)$  such that its
Legendre-Fourier series diverges a.e.. After that, Meaney {\bf[10]} extends the
result to $p=4/3$ and to Jacobi polynomials (see also Badkov {\bf[2]}). In these
results, the asymptotic behavior of the polynomials is used together with
category arguments.

Our purpose in this paper is two fold:

First we obtain results about boundedness of the maximal operators in
weighted $L^p$ spaces, which extend  Badkov's results in the case of Jacobi
polynomials and  are new in the case of Bessel functions (even in the
unweighted case). The procedure is simple and basicly consists of showing that
a general result of Gilbert on transplantation remains true when one considers
weights in the Muckenhoupt $A_p$-classes.

Second we note that  a lower bound on the $L^p$-norm of the polynomials 
is only needed  in order to obtain a result on divergence a.e.. In this
case, a powerful result established by M\'at\'e, Nevai and Totik {\bf[9]}, which is
valid for a very general class of polynomials on $[-1,1]$ allows us to prove a
general divergence theorem which includes the Badkov and Meaney
results. Moreover, we are able to prove an analogous to  M\'at\'e, Nevai and Totik's
result in the context of Bessel functions which allows us to obtain divergence
results in this case.

We shall use the following notations: 

Let $\alpha , \beta >-1$ . We shall denote by $P^{(\alpha , \beta )}_n$ the
Jacobi polynomials on $[-1,1]$ with respect to the measure $d\mu^{(\alpha ,
\beta )} = (1-x)^{\alpha} (1+x)^{\beta} dx$ (see {\bf[15]}).
Let $J_{\alpha}$ be the Bessel function or order $\alpha >-1$, and  
$j_n^{\alpha} (x) = \surd 2 |J_{\alpha+1}(\alpha_n)|^{-1}\por J_{\alpha}
(\alpha_nx)$ (where ${\alpha_n}$ is the increasing sequence of the zeroes of
$J_{\alpha}$)  the Bessel system of order $\alpha$, which is orthonormal
and complete in $L^2 ([0,1];xdx)$.

The corresponding maximal operators will be denoted by $S_{\alpha ,\beta }^*$
and $S_{\alpha }^*$ , respectively.

$A_p([a,b])$ (where $1 < p < \infty$) will stand for the  Muckenhoupt classes
(see {\bf[11]}) consisting of those nonnegative functions $\omega$ on $[a,b]$
such that
$$
   (|I|^{-1} \int_I \omega ) (|I|^{-1} \int_I \omega^{-p'/p})^{p/p'}
   \leq C \qquad {\rm  for\  every\  interval \quad} I\subseteq [a,b],
$$
where $p+p'=p\por p'$.
\medskip
\noindent {\bf \S 1. Extension to $A_p$ classes of a Gilbert result.}
\medskip
\noindent Let $(u_n)_{n\geq 0}$ a sequence of functions on $[0,\pi]$ verifying:

\item{(1.1)} There exists a constant $A$ such that $\displaystyle\sup_{0<x<\pi}
|u_n(x)| \leq A \por , \por \forall n$.

\item{(1.2)} There exist functions $X_1,X_2,X_3,X_4$ in $L^{\infty} ((0,\pi /2))$
such that
$$
   u_n(x) = X_1(x)\cos nx + X_2(x)\sin nx +
   (nx)^{-1}(X_3(x)\cos nx+X_4(x)\sin nx) + O((nx)^{-2})
$$
uniformly for $x \in (1/n,\pi /2)$.

\item{(1.3)} There exist functions $X'_1,X'_2,X'_3,X'_4$ in $L^{\infty}
((0,\pi/2))$ such that
$$
   \eqalign{
   \Delta u_n(x) &= x(X'_1(x)\cos nx+X'_2(x)\sin nx) \cr
   &+ n^{-1}(X'_3(x)\cos nx + X'_4(x)\sin nx) + O(n^{-2}x^{-1})}
$$
uniformly for $x \in (1/n,\pi /2)$, where $\Delta u_n=u_n-u_{n+1}$.

\item{(1.4)} There exists a function $X$ in $L^{\infty} ((0,\pi /2))$ such that
$$
   \Delta u_n(x) =n^{-1}X(x)u_n(x) + (n^{-2}+x)O(1)
$$
uniformly for $x \in (0,1/n)$.

\item{(1.5)} There exists a sequence $(U_n)$ satisfying (1.1)-(1.4) and such that
$$
   u_n(\pi -x)=(-1)^nU_n(x)+O(n^{-2})
$$
uniformly for $x \in (0,\pi /2)$.

Let $(r_{n,k})$ be a double sequence such that:

\item{(1.6)} $\sum_{n=0}^\infty |r_{n,k}| < \infty \por , \por \forall k$.

\item{(1.7)} There exists a constant $B$ such that $|r_{n,k}|\leq B\  \forall
n,k$.

\item{(1.8)} There exists a constant $B$ such that $ \sum_{n=0}^\infty
|\Delta (r_{n,k})| <B \ , \  \forall k$ , where $\Delta
(r_{n,k})=r_{n+1,k}-r_{n,k}$.

If $(u_n),(v_n)$ satisfy (1.1)-(1.5) and $(r_{n,k})$ verifies (1.6)-(1.8) then
we define the kernels:
$$
   K_k(x,t) = \sum_{n=0}^\infty r_{n,k}\por e^{inx}\por e^{-int} \quad,
   \quad L_k(x,t) = \sum_{n=0}^\infty r_{n,k}u_n(x)v_n(t).
$$

In this situation the Gilbert result is (see {\bf [5]}):

\proclaim {Theorem}. Let $1<p<\infty$ . If the operator $f\longrightarrow
\displaystyle\sup_k |\int_{-\pi} ^\pi f(t)K_k(x,t)dt|$ is of weak type $(p,p)$
or strong type $(p,p)$ in $L^p((-\pi ,\pi))$, then the same is true for the
operator $\phi \longrightarrow \displaystyle\sup_k |\int_{-\pi} ^\pi
\phi(t)L_k(x,t)dt|$ in $L^p((0,\pi))$.

Examples of $(u_n)$ satisfying (1.1)-(1.5) are (see {\bf [5]}):
$$
   u_n^{\alpha ,\beta }(\theta) =
   A_n^{\alpha ,\beta }P_n^{\alpha ,\beta }
   (\cos \theta )(1-\cos \theta )^{\alpha /2 +1/4}
   (1+\cos \theta )^{\beta /2 +1/4}\ ,
   \qquad \alpha, \beta \geq -1/2
$$
(where $A_n^{\alpha ,\beta }$ is a normalization constant), or
$$
   u_n^\alpha (\theta )=\pi^{-1}\theta^{1/2}j_n^\alpha (\pi^{-1}\theta )\ ,
   \qquad \alpha  \geq -1/2.
$$

If we call $T_{\alpha ,\beta }^*$ and $T_{\alpha }^*$ to the corresponding
maximal operators, then we have the  relations: 
$$
   S_{\alpha ,\beta}^*f(\cos \theta ) = T_{\alpha ,\beta }^*g(\theta )
  (1-\cos \theta )^{-\alpha /2 -1/4}(1+\cos \theta )^{-\beta /2 -1/4}
  \leqno (1.9)
$$
where $g(\theta)=f(\cos \theta )(1-\cos \theta )^{\alpha /2
+1/4}(1+\cos \theta )^{\beta /2 +1/4}$, and
$$
   S_{\alpha }^*f(\pi^{-1}\theta ) =T_{\alpha }^*g(\theta )\theta^{-1/2}
   \leqno (1.10)
$$
where $g(\theta)=f(\pi^{-1}\theta )\por \theta^{-1/2}$.

It is clear from  Gilbert's theorem that it is possible to obtain boundedness
for the maximal operators $T_{\alpha ,\beta }^*$ and $T_{\alpha }^*$ by
comparing them with the usual Carleson-Hunt maximal operator (see {\bf [4]},
{\bf [6]})  associated to the trigonometric system. However, formulae (1.9) and
(1.10) indicate that we shall need a "weighted" Gilbert theorem if we want to
obtain some boundedness for the maximal operators $S_{\alpha ,\beta }^*$ and
$S_{\alpha }^*$.

\proclaim {Theorem 1}. Let $1<p<\infty$ . Let $\omega \in A_p((-\pi , \pi ))$ an even weight. Let
$(u_n)\ ,\por (v_n)$ and $(r_{n,k})$ be sequences under the hypothesis of 
Gilbert's theorem. If the operator $f\longrightarrow
A^*f=\displaystyle\sup_k |\int_{-\pi} ^\pi f(t)K_k(x,t)dt|$ is of weak  or   
strong type $(p,p)$ in $L^p((-\pi ,\pi),\omega )$, then the same is true for the
operator $\phi \longrightarrow B^*f=\displaystyle\sup_k |\int_{-\pi} ^\pi
\phi(t)L_k(x,t)dt|$ in $L^p((0,\pi),\omega )$.

\proof The Gilbert proof is based on the control of the operator $B^*$
by operators close to the operator $A^*$ (since $(u_n)\ ,\por (v_n)$
may be expressed in terms of sines and cosines) together with another
one which can be essentially controlled by the Hardy-Littlewood maximal
function. Now, the proof consists of observing that all the operators
taking part in the proof are bounded with $A_p$ weights because of the
hypothesis and Muckenhoupt's result related to the boundedness of
the  Hardy-Littlewood maximal function with  $A_p$ weights (see {\bf [11]}).
Details can be found in ({\bf [13]}, Chapter {\bf IV}).

When $r_{n,k}=1$ if $n\leq k$ and $0$ otherwise, and
$(u_n)=(v_n)$ is the sequence $u_n^{\alpha ,\beta }$ (or $u_n^\alpha$), then 
$A^*$ is the operator 
$$
   f\longrightarrow S^*(f^\flat )(x) =
   \displaystyle\sup_k |\sum_{n=0} ^k \hat f(n)e^{inx}|,
$$
where $S^*$ is the Carleson maximal operator and
$f^\flat=\hat f(0)+{1\over2}(f+\tilde f)$ , $\tilde f$ being  the conjugate
function of $f$.

Both operators $S^*$ and $f^\flat$  are bounded in $L^p((-\pi ,\pi),\omega )$
when $\omega$ is an $A_p$ weight, due to  Hunt, Muckenhoupt and Wheeden's
theorem (see {\bf [7]}) and   Hunt and Young's theorem (see {\bf [8]}). On the
other hand, the corresponding operator $B^*$ is the maximal operator
$T_{\alpha ,\beta }^*$ (or $T_{\alpha }^*$ ).

Some applications of Theorem 1 are:

\proclaim {Theorem 2}. 
\item{{\bf (i)}} Let $\alpha,\beta\geq -1/2$  and $1<p<\infty$. Let
$u$ be a weight on $[-1,1]$. If
$$
   \omega (t) = u(\cos t) (1-\cos t)^{(2-p)(2\alpha +1)/4}
   (1+\cos t)^{(2-p)(2\beta +1)/4} \in A_p((0,\pi))
$$
then $S_{\alpha ,\beta }^*$ is bounded in $L^p(u\por d\mu^{(\alpha ,
\beta )})$.
\item{{\bf (ii)}} Let $\alpha \geq -1/2$  and $1<p<\infty$. Let $u$ be a weight
on $[0,1]$.  If $u(x)\por x^{1-p/2} \in A_p((0,1))$  then $S_{\alpha }^*$ is
bounded in $L^p(u(x)\por xdx)$.


\proof {\bf (i)} By (1.9), $S_{\alpha ,\beta }^*$ is bounded on  $L^p(u\por
d\mu^{(\alpha , \beta )})$ if and only if $T_{\alpha ,\beta }^*$ is bounded on 
$L^p(\omega)$ and this follows from the abovementioned remarks. 

{\bf (ii)} It is analogous using (1.10) instead of (1.9).

A particular case of part (i) in Theorem 2 is

\proclaim Corollary 1 (Badkov {\bf [2]}). Let $\alpha, \beta \geq -1/2$, $1 < p
< \infty$, $u(x)=(1-x)^{ap}(1+x)^{bp}$ . If the conditions:
$$
   \eqalign{
   & a + (\alpha+1)(1/p-1/2) < 1/4, \quad b + (\beta+1)(1/p-1/2) < 1/4 \cr
   & a + (\alpha+1)(1/p-1/2) > -1/4 ,\quad b + (\beta+1)(1/p-1/2) > -1/4
   }
$$
are fullfilled, then 
$S_{\alpha ,\beta }^*$ is bounded on $L^p(u\por d\mu^{(\alpha , \beta )})$.

\proof We need only to check that 
$$
   (1-\cos t)^{ap+\{(2-p)(2\alpha +1)\}/4}
   (1+\cos t)^{bp+\{(2-p)(2\beta +1)\}/4} \in A_p((0,\pi))
$$
but this is equivalent to
$$
   t^{2ap+\{(2-p)(2\alpha +1)\}/2}
   (\pi - t)^{2bp+\{(2-p)(2\beta +1)\}/2} \in A_p((0,\pi))
$$
and a simple exercise of integration shows that this is true taking into account
the hypothesis.

\noindent {\bf Note.}  Considering different systems $(u_n),(v_n)$, theorem 1
also allows us to obtain transplantation's results of the following type:

{\sl ``Let $d\mu,d\nu$ be two different measures of Jacobi type and
let $P_n,Q_n$ the corresponding orthonormal Jacobi polynomials. Then, for a
certain range of $p's$, for each $f\in L^p(d\nu)$, there exists a $g\in
L^p(d\mu)$ such that 
$$
\int_{-1}^1 g(x)P_n(x)d\mu(x)=\int_{-1}^1 f(x)Q_n(x)d\nu(x)\qquad \forall
n\geq 0
$$
and \quad$||g||_{L^p(d\mu)} \leq C\por ||f||_{L^p(d\nu)}$\ \  where the constant
$C$ is independent of $f$".}

Theorems of this type have been studied by Askey {\bf [1]} and Muckenhoupt {\bf
[12]}, not for the polynomials on $[-1,1]$ but for the corresponding
orthonormal functions in $(0,\pi)$.

\medskip
\noindent {\bf \S 2.  Divergence almost everywhere.}
\medskip
\noindent 
First, we are going to study the case of orthonormal polynomials in $[-1,1]$.

Assume $d\mu$ is a positive Borel measure on $[-1,1]$ such that $\mu'> 0$ a.e., 
and let $\{P_n(x)\}_{n=0}^\infty$ denote the corresponding orthonormal
polynomials.

The following result is due to M\'at\'e, Nevai and Totik {\bf [9]}:

\proclaim {Lemma 1}. Suppose $1\leq r < \infty$. If $g$ is a Lebesgue measurable
function in $[-1,1]$, then
$$
   (\int_{-1}^1 |g(x)(1-x^2)^{-1/4}\mu'(x)^{-1/2}|^r dx)^{1/r} \leq
   M\por \liminf_{n\to \infty} (\int_{-1}^1 |g(x)P_n(x)|^r dx)^{1/r}
$$
where $M= \pi^{1/2}2^{max\{1/r-1/2,0\}}$.

By using this lemma we obtain:

\proclaim {Theorem 3}. Let $\omega$ be a Lebesgue
measurable positive function in $[-1,1]$  and $1<p\leq q < \infty$,
$(1/p+1/q=1)$. If 
$$
   \omega(x)^{-1} \not\in L^q(d\mu) \qquad or \qquad
   \mu'(x)^{1-q/2}(1-x^2)^{-q/4} \omega(x){-q} \not\in L^1(dx), \leqno (2.1)
$$
then there exists $f\in L^p(\omega^pd\mu)$ such that $S_nf$ either does not
exist or diverges almost everywhere in $[-1,1]$.

\proof Suppose that $c_0(f) = \int_{-1}^1 f(x) P_0(x) d\mu(x)$ exists for each
$f \in L^p(\omega^pd\mu)$. By duality, this implies $P_0 \in
L^q(\omega^{-q}d\mu)$, i.e., $\omega^{-1} \in  L^q(d\mu)$. If  $\omega^{-1} \in 
L^q(d\mu)$ , then $S_nf$ exists for every $n=0,1,2,...$ and every $f\in
L^p(\omega^pd\mu)$ , and so the operators
$$
   c_n:f\in L^p(\omega^pd\mu) \longrightarrow
   c_n(f)=\int_{-1}^1f(x)P_n(x)d\mu(x)
$$
are continuous by using Holder's inequality.
Moreover, $\Vert c_n \Vert = \Vert P_n\omega^{-1}\Vert_{L^q(d\mu)}$ and so, by
(2.1) and lemma 1, it is clear that
$\displaystyle \sup_n \Vert P_n\omega^{-1} \Vert_{L^q(d\mu)} = \infty$. Then,
the  Banach-Steinhaus theorem shows that there exists $f \in L^p(\omega^pd\mu)$
such that $\displaystyle \sup_n |c_n(f)| = \infty$.

On the other hand, if $S_nf(x)=\sum_{k=0}^n c_k(f) P_k(x)$ converges in a
set $E$ of Lebesgue measure $|E|>0$ and $\displaystyle \sup_n |c_n(f)| =
\infty$, then there exists a subsequence $P_{n_k}$ with
$$
   \lim_k P_{n_k}(x) = 0 \qquad {\rm a.e.\ in\ } E.
$$
If we take $\delta>0$ and $\delta<|E|$, by using Egoroff's theorem, there exists
$D\subset E$ such that $|D|=\delta$ and $\displaystyle \lim_k P_{n_k}(x) = 0$
uniformly in $D$. Therefore
$$
   \lim_k \int_D |P_{n_k}| d\mu = 0.
$$
But, by lemma 1 with $r=1$ and $g(x) = \mu'(x) \chi_D(x)$, we have
$$
   \lim_k \int_D |P_{n_k}| d\mu \geq
   M^{-1} \int_D (1-x^2)^{-1/4} \mu'(x)^{1/2} dx > 0,
$$
which is a contradiction.

Next, we show a similar result for Bessel systems. First of all, we need to
establish an analogous result to lemma 1. For that, we shall use the following
estimates for Bessel functions
$$
   J_\alpha(z) = z^\alpha2^{-\alpha}\Gamma(\alpha+1)^{-1} + O(z^{\alpha+2}),
   \quad z \longrightarrow 0 \leqno (2.2)
$$
$$
   J_\alpha(z) = \sqrt 2 (\pi z)^{-1/2}
   \cos (z-{\alpha \pi \over 2}-{ \pi \over 4}) + O(z^{-3/2})
   \quad z \longrightarrow \infty \leqno (2.3)
$$
$$
   \int_0^1 J_\alpha(\alpha_nx)^2 x\por dx \approx {1 \over \pi\alpha_n},
   \quad n \longrightarrow \infty \leqno (2.4)
$$
and the Fejer lemma: if $f,g$ are nonnegative functions, $g$
is continuous with period $2\pi$ and $f\in L^1((0,2\pi))$, then
$$
   \lim_{\lambda \to \infty} {1 \over 2 \pi}
   \int_0^{2\pi} f(t) g(\lambda t) dt = \hat g (0) \por \hat f (0) \leqno (2.5)
$$
where $\hat h (x)$ denotes the Fourier transform of $h$.

\proclaim {Lemma 2}. Let $\alpha>-1$. Let $h(x)$ be a Lebesgue measurable
nonnegative function on $[-1,1]$ and $1<p<\infty$. Then
$$
   \lim_{n \to \infty} \int_0^1 |j_n^{\alpha}(x)|^p h(x) dx \geq
   M \por \int_0^1 h(x) \por x^{-p/2} \por dx \leqno  (2.6)
$$
where
$$
   M = {2^{p/2} \over 2\pi} \int_0^{2\pi} |\cos t|^p \por dt.
$$

\proof If there exists no $\delta \in (0,1)$ such that $h(x)x^{\alpha p} \not
\in L^1 (0,\delta)$ then (2.6) is trivial since (2.2) implies that all the
integrals on the left in (2.6) are equal to $\infty$. Then, we may suppose that 
$h(x) x^{\alpha p} \in L^1 (0,\delta)$ for some $\delta \in (0,1)$. By (2.4) it
is enough to prove
$$
   \lim_{n \to \infty}
   \int_0^1 |(\pi\alpha_n)^{1/2}J_\alpha(\alpha_nx)|^p h(x) dx =
   M \por \int_0^1 h(x) \por x^{-p/2} \por dx. \leqno  (2.7)
$$
Suppose that $h(x)x^{-p/2}\in L^1((0,1))$. If we prove 
$$
   \lim_{n \to \infty} \int_0^1
   | (\alpha_n x)^{1/2} J_\alpha (\alpha_n x) -
   (2/\pi)^{1/2} \cos (\alpha_n x - {\alpha \pi \over 2}-{\pi \over 4}) |^p
   h(x) x^{-p/2} dx = 0 \leqno (2.8)
$$
then (2.7) is reduced to show that
$$
   \lim_{n \to \infty} \int_0^1
   | (2/x)^{1/2} \cos (\alpha_n x - {\alpha \pi \over 2} - { \pi \over 4}) |^p
   h(x) dx =
   M  \por \int_0^1 h(x) \por x^{-p/2} \por dx. \leqno (2.9)
$$
Now, (2.9) holds by using (2.5) with 
$$
   t = 2 \pi x,
   \ g(t) = |2^{1/2} \cos (t - {\alpha\pi \over 2} - { \pi\over 4})|^p,
   \ f(t) = h({t \over 2 \pi}) \Big({t \over 2 \pi} \Big)^{-p/2},
   \ \lambda = {\alpha_n \over 2 \pi}.
$$
Thus, let us prove (2.8) under conditions 
$$
   h(x) x^{\alpha p}\in L^1((0,\delta)) \quad {\rm and} \quad
   h(x)x^{- p/2} \in L^1((0,1). \leqno (2.10)
$$
Take $n$ large enough so that $\alpha_n>1$. In view of (2.4), if $\alpha_n x\geq
1$ then
$$
   | (\alpha_n x)^{1/2} J_\alpha (\alpha_n x) -
   (2/\pi)^{1/2} \cos (\alpha_n x - {\alpha \pi \over 2} - {\pi \over 4}) |
   \leq C \por (\alpha_n x)^{1/2} (\alpha_n x)^{-3/2} < C_1
$$
and, according to (2.2), if $\alpha_n x\leq 1$ it follows
$$
   | (\alpha_n x)^{1/2} J_\alpha (\alpha_n x) -
   (2/\pi)^{1/2} \cos (\alpha_n x - {\alpha \pi \over 2} - {\pi \over 4}) |
   \leq C \por ((\alpha_n x)^{\alpha + 1/2} + 1)
   \leq  C_1 \por (x^{\alpha + 1/2} + 1)
$$
(in the last inequality we use $\alpha_nx \leq 1$, if  $\alpha+1/2\geq 0$,
and $\alpha_nx \geq x$, if $\alpha+1/2\leq 0$,  ).

Now, by (2.10), we can apply the dominate convergence theorem
and we obtain (2.8). 

Finally, in case that $h(x) x^{-p/2} \not \in L^1((0,1))$, let
$\{K_j\}_{j=1}^\infty$ be the sequence of increasing measurable sets
$K_j=\{x\in (0,1) : h(x) x^{-p/2} \leq j\}$, and denote by $h_j$ the function
$h_j=h$ on  $K_j$ and $0$ elsewhere. By applying (2.7) for each  $h_j$ and then 
the monotone convergence theorem, (2.7)  also holds in this case, and the proof
of the lemma is finished.

From lemma 2 and proceeding exactly as in the proof of theorem 3, we obtain:

\proclaim {Theorem 4}. Let $\alpha>-1$ be and suppose $1<p\leq q < \infty$, 
$(1/p+1/q=1)$. Let $\omega$ be a Lebesgue measurable nonnegative function on
$[0,1]$. If
$$
   \omega(x)^{-q}x^{\alpha q+1} \not\in L^1(dx) \quad or \quad 
   \omega(x)^{-q}x^{1 - q/2} \not\in L^1(dx)
$$
then there exists $f \in L^p(\omega(x)^p x dx)$ whose Fourier series with
respect to the system $\{j_n^\alpha(x)\}$ either does not exist or diverges
almost everywhere on $[0,1]$.

\noindent {\bf Notes.}

\item {{\bf (i)}} If $\alpha \geq -1/2$, the proof of lemma 2 shows that the
sign ``$\geq$'' can be replaced by ``$=$''.

\item {{\bf (ii)}} Theorems 3 and 4 contain as particular cases  the known cases
of generalized Jacobi polynomials ({\bf [2]}, {\bf [10]}, {\bf [14]}) or Bessel
systems ({\bf [3]}).

\item {{\bf (iii)}} An analogous result to Theorem 4 can be obtained for Dini
systems in $[0,1]$.

\medskip

{\bf  References.} \frenchspacing

\medskip

\item {{\bf [1]}} R. Askey, {\sl A transplantation theorem for Jacobi series}, 
Illinois J. Math. {\bf 13 } (1969), 583-590.

\item {{\bf [2]}} V. M. Badkov, {\sl Convergence in the mean and almost
everywhere of Fourier series in polynomials orthogonal on an interval},  Math.
USSR Sb. {\bf 24} (1974), 223-256.

\item {{\bf [3]}} A. I. Benedek and R. Panzone, {\sl Pointwise convergence of
series of Bessel functions}, Rev. Un. Mat. Arg. {\bf 26} (1972), 167-186.

\item {{\bf [4]}} L. Carleson, {\sl On convergence and growth of partial sums of
Fourier series}, Acta Math. {\bf 116} (1966), 135-157.

\item {{\bf [5]}} J. E. Gilbert, {\sl Maximal theorems for some orthogonal
series. I}, Trans. Amer. Math. Soc. {\bf 145} (1969), 495-515.

\item {{\bf [6]}} R. Hunt, {\sl On the convergence of Fourier series}, Proc.
Conf. Orthogonal Expansions and Continuous Analogues, 235-255. Southern Illinois
Univ. Press, Carbondale, Ill., 1968.

\item {{\bf [7]}} R. Hunt, B. Muckenhoupt and R. Wheeden, {\sl Weighted norm
inequalities for the conjugate function and Hilbert transform}, Trans. Amer.
Math. Soc. {\bf 176} (1973), 227-251.

\item {{\bf [8]}} R. Hunt and W. S. Young, {\sl A weighted norm inequality for
Fourier series}, Bull. Amer. Math. Soc. {\bf 80} (1974), 274-277.

\item {{\bf [9]}} A. M\'at\'e, P. Nevai and V. Totik, {\sl Necessary conditions
for weighted mean convergence of Fourier series in orthogonal polynomials}, J.
Approx. Theory {\bf 46} (1986), 314-322.

\item {{\bf [10]}} C. Meaney, {\sl Divergent Jacobi polynomial series}, Proc.
Amer. Math. Soc. {\bf 87} (1983), 459-462.

\item {{\bf [11]}} B. Muckenhoupt, {\sl Weighted norm inequalities for the Hardy
maximal function}, Trans. Amer. Math. Soc. {\bf 165} (1972), 207-226.

\item {{\bf [12]}} B. Muckenhoupt, {\sl ``Transplantation theorems and
multiplier theorems for Jacobi series"}, Memoirs of the Amer. Math. Soc., vol.
64, n� 356, Providence, R. I., 1986.

\item {{\bf [13]}} M. P\'erez, {\sl Series de Fourier respecto de sistemas
ortogonales: estudio de la convergencia en espacios de Lebesgue y de
Lorentz}, Doctoral Dissertation. Sem. Mat. Garc\'\i a de Galdeano,
Zaragoza, 1989.

\item {{\bf [14]}} H. Pollard, {\sl The convergence almost everywhere of Legendre
series}, Proc. Amer. Math. Soc. {\bf 35} (1972), 442-444.

\item {{\bf [15]}} G. Szeg\H o, {\sl \lq\lq Orthogonal polynomials"}, 3rd. ed.,
Amer. Math. Soc., Providence, R. I., 1967.

\item {{\bf [16]}} E. C. Titchmarsh, {\sl\lq\lq Eigenfunction expansions
associated with second-order differential equations"}, Vol. 1, 2nd. ed., Oxford
Univ. Press, Oxford, 1962.

\bigskip

\noindent {\rmdiez Jos\'e J. Guadalupe, Mario P\'erez and Francisco
J. Ruiz. Departamento de Matem\'aticas. Facultad de Ciencias. Universidad de
Zaragoza, 50009 Zaragoza, SPAIN.}

\noindent {\rmdiez Juan L. Varona. Departamento de Matem\'atica
Aplicada. Colegio Universitario de La Rioja, 26001 Logro\~no, SPAIN.}





 \bye
